\section{Privacy、Encryption 与最小必要信号}

安全系统会处理高度敏感内容与账户信息，隐私不是与安全对立的“外部限制”，而是架构目标。\term{end-to-end encryption}（端到端加密，E2EE）意味着消息内容只有通信端点持有密钥，传输和服务端无法读取明文；它保护用户免受平台、攻击者和中间方的内容窥探，也限制服务端 content scanning。负责任设计应明确哪些信号仍可用、谁可访问、如何防止扩大监控。

\subsection{Content、metadata、endpoint 与 governance 四层}

本节不寻找一个“绕过加密”的万能方案，而是分层讨论。\term{metadata}（元数据）是内容之外的通信或系统属性，例如时间、账户、设备、连接关系和服务事件；它可能具有强隐私敏感度，不能因为不是正文就无限收集。Endpoint controls 位于用户设备或应用边界，network signals 描述协调行为，governance 则约束访问、legal process、appeal 与 oversight。

\lecturefigure{11-privacy-layers.png}{隐私与儿童安全需要 content、metadata、endpoint 和 governance 的分层控制。}{本地字幕 28:40--31:10；NIST/平台治理原则；概念重绘。}

\begin{knowledgebox}{读图：每种信号都写 purpose limitation}
Content signal 可能最直接但侵入性最高；metadata 能显示异常关系却容易误伤普通社交行为；endpoint intervention 可以更早提示用户，却需要防止设备级监控扩张；governance 决定谁可查询、何时删除、怎样申诉。对每层都应写明 purpose、合法依据、最小字段、retention、access role、audit 和禁止的 secondary use。
\end{knowledgebox}

\begin{warningbox}{“隐私与安全都重要”不是设计答案}
真正的工程答案需要 threat model 和边界：是否改变端到端机密性，是否引入可被其他目的复用的扫描能力，误报由谁纠正，未成年人和弱势群体是否承受额外风险，政府或内部人员能否扩大查询。任何方案都应接受独立安全评估、权利影响评估和透明治理，而不是只靠供应商承诺。
\end{warningbox}

\teachervoice{Cordua 承认 encryption 带来真实限制，同时强调平台仍要寻找合法、隐私保护的安全路径。课堂最值得保留的不是某个单一政策立场，而是“不要假装 trade-off 不存在，也不要因此停止做系统设计”。}

\subsection{本章小结}

隐私与安全都需要最小权限、purpose limitation 和可审计治理。E2EE 保护内容机密性并改变服务端可见性；metadata、endpoint 或 network signal 只能作为分层风险信号，不能自动成为广泛监控的正当化理由。

